\scanpage{245}
% paper-clip impression at upper left
donc pour $n$ grand \inserted{$(\geq n_0)$}, et $m\geq n$,
\[
\check H^{\unsure{*|\times}}(k'/k,M_m)\overset{\sim}{\longrightarrow}\check H^*(k'/k,M_n)
\]

\deleted{Or on a un homomorphisme}

\deleted{D'où facilement} Donc si
$\unsure{\xi|\zeta}_{\unsure{m|n}}\in\underline{\operatorname{Pic}}_{X_m/k}(k)$ est tel que son image
$\xi_n$ dans $\underline{\operatorname{Pic}}_{X_n/k}(k)$ provient de $\underline{\operatorname{Pic}}(X_n)$, alors
il provient lui-$=$ de $\underline{\operatorname{Pic}}(X_m)$. Ceci, et
le fait que $(\underline{\operatorname{Pic}}_{X_n/k}(k))$ satisfait à
M.L. (A.R., montre que il en est
de $=$ de $(\operatorname{Pic}(X_n))_n$, de façon
précise si $n\geq n_0$, \deleted{\unsure{alors il y a|tout élément}}
et si \deleted{\unsure{$m$ est}}\inserted{$n'$} est tel que $\operatorname{Im}(\underline{\operatorname{Pic}}_{X_{\deleted{\unsure{m}}n'}/k}(k)\unsure{-|\to}\underline{\operatorname{Pic}}_{X_n/k}(k)\unsure{)}$
soit \unsure{minimal|minimale|universel}, alors $\operatorname{Im}(\operatorname{Pic}(X_{\deleted{\unsure{m}}n'})\to\operatorname{Pic}(X_n))$
est \unsure{minimal|minimale|universel}. OK.

\diagram{\begin{tikzcd}[ampersand replacement=\&,column sep=small,row sep=small]
0\arrow[r,dash] \& \operatorname{Pic}(X_n)\arrow[d,phantom,"\rotatebox{-45}{$\ni$}"]\arrow[r,dash] \& \underline{\operatorname{Pic}}_{X_n/k}(k)\arrow[d,phantom,"\rotatebox{-45}{$\ni$}"]\arrow[r,dash] \& \deleted{\check H^2(k'/k,M_n)}\\
\& \xi_n\arrow[r,decorate,decoration={snake,amplitude=.35mm,segment length=2mm,post length=1mm}] \& \bar\xi_n \& \\
0\arrow[r,dash] \& \operatorname{Pic}(X_{n'})\arrow[uu]\arrow[r]\arrow[d,phantom,"\rotatebox{-45}{$\ni$}"] \& \underline{\operatorname{Pic}}_{X_{n'}/k}(k)\arrow[uu] \& \\
\& \xi_{n'}\arrow[uu,decorate,decoration={snake,amplitude=.35mm,segment length=2mm,post length=1mm}] \& {} \& \\
0\arrow[r,dash] \& \operatorname{Pic}(X_{n''})\arrow[uu]\arrow[r,dash]\arrow[d,phantom,"\rotatebox{-45}{$\ni$}"] \& \underline{\operatorname{Pic}}_{X_{\unsure{k''|n''}}/k}(k)\arrow[uu]\arrow[d,phantom,"\rotatebox{-45}{$\ni$}"] \& \\
\& \deleted{\unsure{\xi_m}}\lambda_m\arrow[r,decorate,decoration={snake,amplitude=.35mm,segment length=2mm,post length=1mm}] \& \unsure{\eta|\zeta}_m\arrow[uuu,decorate,decoration={snake,amplitude=.35mm,segment length=2mm,post length=1mm},bend right=20] \&
\end{tikzcd}}

\uline{Corollaire} Pour que
$\mathfrak X/\unsure{\mathcal Y|\mathfrak Y|Y}$ soit \uline{projectif},
il f. et s. que pour
tt $n$, $X_n/Y_n$ \inserted{le} soit.
% archive sheet number: 244
